\section{Introduction}
\label{sec:intro}

Object detection at real-time rates is dominated by two paradigms: NMS-based one-stage detectors (the YOLO family~\cite{redmon2016yolo,wang2023yolov7}) and end-to-end transformer detectors that learn a permutation-invariant set prediction (DETR~\cite{carion2020detr} and its descendants~\cite{zhu2021deformable,liu2022dabdetr,zhang2023dino}). RT-DETR~\cite{lyu2024rtdetr} closed the speed gap between these paradigms by introducing a hybrid encoder that decouples intra-scale interaction (Attention-based In-scale Feature Interaction, AIFI) from cross-scale fusion (CNN-based Cross-scale Feature Fusion, CCFF), and demonstrating that a transformer decoder with bipartite matching can match or exceed YOLO accuracy at comparable FPS, with the operational advantage of NMS-free inference.

A persistent weakness of the DETR family --- and RT-DETR is no exception --- is small-object detection. On COCO~\cite{lin2014coco} val2017, RT-DETR-R50 reports $\mathrm{AP}_S = 34.7$ at $640\!\times\!640$, well below its $\mathrm{AP}_L = 67.7$ and a couple of points below comparable CNN-based detectors that benefit from explicit feature pyramids targeting stride-4/8 features~\cite{lin2017fpn}. The gap is consequential: small-object performance is the binding constraint in many practical applications (aerial, surveillance, traffic).

The hypothesis we explore in this work is that the encoder memory --- a flat sequence of multi-level tokens cross-attended over by every decoder layer --- is computed once and never refined as the decoder accumulates information about object hypotheses. The information flow is strictly forward. We ask: \emph{would feeding back early decoder predictions to refine the encoder memory help small-object detection?} Our thinking is that small objects benefit disproportionately from having every available cue used at every layer, because they have so few pixels to begin with.

\paragraph{Contributions.}
\begin{enumerate}
    \item We introduce a \emph{Decoder-to-Encoder Feedback} module: after the first decoder layer produces preliminary predictions, its query output is used as keys/values in a cross-attention pass that refines the encoder memory. Later decoder layers then attend over the refined memory.
    \item We empirically show that a naive implementation (v1) is statistically silent at inference: a same-checkpoint ablation produces $\Delta \mathrm{AP}_S = 0.0$ between feedback ON and OFF. We trace the cause to gate collapse: the learnable gate $g = \sigma(\alpha)$ that mixes the refinement into the memory decays to ${\approx}0.12$ during training, effectively silencing the mechanism.
    \item We propose two structural fixes that, together, turn the silent mechanism into a measurable one: (a) a \emph{gate floor} reparameterization $g_{\mathrm{eff}} = \mathrm{floor} + (1{-}\mathrm{floor})\cdot\sigma(\alpha)$, which makes the gate physically incapable of collapsing, and (b) a \emph{level mask} that restricts cross-attention to the P2 (stride 4) and P3 (stride 8) features where small objects live, focusing the gradient signal on the metric we want to move.
    \item In a same-checkpoint ablation, v2 yields $\Delta\mathrm{AP}_S = +0.99$ ($33.26 \to 34.25$), with $\Delta\mathrm{AP} = +0.57$, $\Delta\mathrm{AP}_M = +0.42$, $\Delta\mathrm{AP}_L = +0.50$. All four metrics improve under feedback ON, with $\mathrm{AP}_S$ improving the most --- exactly as the hypothesis predicted. Critically, v2 adds zero new parameters relative to v1; the changes are purely a reparameterization and a compute-time mask.
\end{enumerate}

We see this as a small contribution methodologically but a useful negative-result-turned-positive: the same architectural ideas can be silenced or made productive depending on a single structural choice in how the gating is parameterized. We document both outcomes in detail because the contrast is the empirical point.
